\section{Roteiro Prático 01 — Retomada e Transição Parte 1: Bootstrap 5 \& Layout Responsivo Estático}

% =====================================================
% OBJETIVOS
% =====================================================

\begin{boxObjetivos}
\textbf{Objetivos da Aula (Parte 1 — Bootstrap 5)}

Ao final desta aula o estudante será capaz de:

\begin{itemize}
\item Conectar a base de HTML5 semântico e CSS3 de \textbf{ProgWeb 1} com a evolução prática de \textbf{ProgWeb 2};
\item Compreender o alinhamento do semestre com o Projeto Pedagógico do Curso (\textbf{PPC} CONSUP 30/2020 — 120h);
\item Localizar e obter a biblioteca no site oficial do \textbf{Bootstrap 5} (\url{https://getbootstrap.com});
\item Configurar o ambiente incluindo os arquivos de estilo CSS e Bootstrap Icons no \texttt{<head>} do documento HTML;
\item Aplicar o sistema de \textbf{Grid Responsivo de 12 Colunas} do Bootstrap 5 (\texttt{.container}, \texttt{.row}, \texttt{.col-md-*});
\item Estilitar elementos visuais da interface (Exemplos 1 a 5: Grid, Formulários, Cards, Badges, Modais e Tabelas Zebradas);
\item Estruturar a interface do \textbf{Layout Estático Completo da Aplicação Web} (Exemplo 5);
\item Acompanhar o material através da apresentação de slides interativa da Parte 1 (\url{https://chameoandre.github.io/CST-Sistemas-para-internet-programacao-para-internet-2/retomada-de-atividades-pi2/parte-1-bootstrap/index.html}).
\end{itemize}
\end{boxObjetivos}

% =====================================================
% INTRODUÇÃO & INCLUSÃO
% =====================================================

\subsection{Introdução, Site Oficial e Inclusão do Bootstrap 5}

Bem-vindos à primeira parte da retomada prática da disciplina de \textbf{Programação para a Internet 2 (ProgWeb 2)}!

Neste primeiro encontro, focamos na construção da interface visual responsiva estática do nosso sistema utilizando o framework \textbf{Bootstrap 5}.

\begin{enumerate}
\item \textbf{Bootstrap 5 (Site Oficial \& CDN):} No portal oficial \url{https://getbootstrap.com} e nas referências do W3Schools, obtemos a documentação de componentes e os links de CDN (\texttt{jsDelivr}) incluídos via tag \texttt{<link rel="stylesheet">} na seção \texttt{<head>};
\item \textbf{Bootstrap Icons:} Inclusão do pacote de ícones vetoriais oficial do Bootstrap para ilustrar botões, cartões e alertas da aplicação;
\item \textbf{Acompanhamento Interativo:} O trabalho é guiado pelos slides da Parte 1 (\texttt{parte-1-bootstrap/index.html}) e pela prática na pasta \texttt{exercicio-retomada/}.
\end{enumerate}

% =====================================================
% PARTE 1
% =====================================================

\subsection{Detalhamento Didático dos Exemplos de Bootstrap 5 (Exemplos 1 a 5)}

Abaixo está o detalhamento dos componentes de layout desenvolvidos nesta etapa:

\begin{itemize}
\item \textbf{Exemplo 1 (Bootstrap 5 — Grid Responsivo \& Containers):} \textit{Desenvolvimento de layout responsivo flexível.} A classe \texttt{.container} centraliza o conteúdo, \texttt{.row} cria a linha de grade e \texttt{.col-12 .col-md-6 .col-lg-4} ajusta a largura para 100\% no celular, 50\% em tablets e 33.3\% em desktops;
\item \textbf{Exemplo 2 (Bootstrap 5 — Formulários \& Inputs Estilizados):} \textit{Construção de campos com visual moderno.} A classe \texttt{.mb-3} adiciona margem inferior, \texttt{.form-label} formata o rótulo e \texttt{.form-control} estiliza o campo com bordas arredondadas e efeito de foco azul;
\item \textbf{Exemplo 3 (Bootstrap 5 — Cards, Botões \& Alertas Visuais):} \textit{Empacotamento de componentes de interface.} Cria um cartão com elevação visual (\texttt{.card .shadow}), exibe um alerta verde de confirmação (\texttt{.alert .alert-success}) e insere um botão estilizado (\texttt{.btn .btn-primary});
\item \textbf{Exemplo 4 (Bootstrap 5 — Tabelas Elegantes \& Utilitários):} \textit{Formatação de tabelas com múltiplos registros.} Aplica \texttt{.table} para espaçamento automático, \texttt{.table-hover} para destacar a linha sob o mouse, \texttt{.table-striped} para alternar as cores das linhas (zebrado visual) e \texttt{.badge} para exibir tags coloridas de status;
\item \textbf{Exemplo 5 (Bootstrap 5 — Layout Estático Completo):} \textit{Layout integrado de cadastro e listagem.} Combina Grid responsivo (\texttt{.row .col-lg-4} e \texttt{.col-lg-8}), Formulário no Card e Tabela Zebrada em uma única interface completa, pronta para receber a interatividade na Parte 2.
\end{itemize}

% =====================================================
% TABELAS "PARA EXPLORAR MAIS"
% =====================================================

\subsection{Guia de Apoio: "Para Explorar Mais no Bootstrap 5"}

Para dar suporte aos alunos na customização do layout, o quadro abaixo resume opções adicionais de classes do Bootstrap 5:

\begin{boxTabela}
\textbf{Opções Adicionais do Bootstrap 5 (Layout, Componentes e Utilitários)}

\begin{tabular}{|l|p{11cm}|}
\hline
\textbf{Código / Classe} & \textbf{Descrição e Efeito Visual} \\ \hline
\texttt{col-12 col-sm-6 col-lg-3} & Ocupa 100\% no celular, 50\% em tablets pequenos e 25\% em desktops wide. \\ \hline
\texttt{form-select / form-control-lg} & Aplica o estilo de menu suspenso ou aumenta o tamanho do campo input. \\ \hline
\texttt{btn-outline-danger} & Botão vazado vermelho que preenche ao passar o ponteiro do mouse. \\ \hline
\texttt{shadow-sm / shadow-lg} & Aplica sombra suave ou pronunciada para destacar cartões da página. \\ \hline
\texttt{d-flex justify-content-between} & Ativa o Flexbox e alinha itens nas extremidades da caixa. \\ \hline
\end{tabular}
\end{boxTabela}

% =====================================================
% TAREFA PRÁTICA
% =====================================================

\subsection{Tarefa Prática de Laboratório (Construção do Layout Estático)}

\begin{boxAtividade}[{Atividade Prática -- Layout Estático com Bootstrap 5}]

Na pasta \texttt{exercicio-retomada/}, os estudantes deverão construir a estrutura estática:

\begin{enumerate}
\item \textbf{Inclusão dos Links no Head:} Abrir o arquivo \texttt{index.html} e incluir os links de CDN do Bootstrap 5 CSS e Bootstrap Icons;
\item \textbf{Montagem da Grade Responsiva:} Criar a estrutura do Grid dividindo a página entre a coluna de cadastro (\texttt{col-md-4}) e a coluna da tabela (\texttt{col-md-8});
\item \textbf{Estilização de Formulários e Tabelas:} Aplicar as classes \texttt{.card}, \texttt{.form-control}, \texttt{.btn-primary} e \texttt{.table-striped} construindo o leiaute profissional completo (Exemplo 5).
\end{enumerate}

\end{boxAtividade}

% =====================================================
% CHECKLIST
% =====================================================

\subsection{Checklist Final da Aula 1}

\begin{boxAtividade}[{Verificação Técnica Parte 1}]

Antes de finalizar o encontro, verifique se seu projeto atende aos critérios de aceitação:

\begin{itemize}
\item [ ] As dependências do Bootstrap 5 CSS e Icons estão corretamente vinculadas no \texttt{<head>};
\item [ ] A página ajusta seu layout de forma fluida ao alterar a largura da janela do navegador;
\item [ ] A interface estática da aplicação (Exemplo 5) está completamente montada e legível.
\end{itemize}

\end{boxAtividade}
